\documentclass[aspectratio=43]{beamer}
\usetheme{Berlin}

\usepackage[english]{babel}
\usecolortheme{dolphin}
\usepackage{graphicx}
\usepackage{dirtree}
\usepackage{listings}
\usepackage[T1]{fontenc}
\usepackage{lmodern}
\usepackage[utf8]{inputenc}
\usepackage{caption}
\usepackage{bbding}
\usepackage{xurl}
\usepackage{scrextend}
\usepackage{minted}
\usepackage{appendixnumberbeamer}
\usepackage[backend=bibtex, style=numeric, maxcitenames=1, isbn=false, url=false, doi=false, eprint=false]{biblatex}

\addbibresource{../../doc/bibliography.bib}

\captionsetup{labelformat=empty}

\beamertemplatenavigationsymbolsempty
\defbeamertemplate*{title page}{customized}[1][]
{
\usebeamerfont{title}\inserttitle\par
\usebeamerfont{subtitle}\usebeamercolor[fg]{subtitle}\insertsubtitle\par
\bigskip
\usebeamerfont{author}\insertauthor\par
\usebeamerfont{institute}\insertinstitute\par
\usebeamerfont{date}\insertdate\par
\usebeamercolor[fg]{titlegraphic}\inserttitlegraphic
}

\hypersetup{unicode}
\hypersetup{breaklinks=true}

\title{3D-Printed Metamaterials}
\subtitle{Prepared for meeting on 9.24.2026}
\author{Havránek Kryštof}
\date{}
\institute{CTU, NTUST}
\setbeamertemplate{sidebar right}{}
\setbeamertemplate{footline}{%
\hfill\textbf{Page \insertframenumber{}. z \inserttotalframenumber} \hspace{0.01cm} \vspace{0.1cm}}
\setbeamerfont{footnote}{size=\tiny}

\begin{document}

\begin{frame}[plain]
  \maketitle
\end{frame}

\clearpage
\setcounter{framenumber}{0}

\begin{frame}[fragile]
	\section{Introduction}
  \frametitle{Outline}
  \begin{itemize}
    \item Topic of the thesis.
    \item High permittivity and conductive material situations.
    \item Prospective topics in leaky-wave antennas.
    \item Prospective topics in reflectarray antennas.
  \end{itemize}
\end{frame}

\begin{frame}[fragile]
  \frametitle{Topic of the thesis}
  \begin{itemize}
    \item "Metamaterial" in the loosest sense, not limited to left handed materials.
		\item Focused only on topics utilizing dielectric, or lossy filaments
		\item Conventional 3DP techniques - FDM, SLA, SLS
		\item Specific topic is still to be decided.
  \end{itemize}
\end{frame}

\begin{frame}[fragile]
  \frametitle{High Permittivity Materials}
	\section{Materials}
  \begin{itemize}
		\item High permittivity enables more applications, or more compact design with lower losses.
		\item FDM
			\begin{itemize}
				\item Zetamix: $\varepsilon_r=10$ with $\tan \delta 10^{-3}$, B2C
				\item TPU: high permittivity, however $\tan \delta \approx 0.03 ~ 0.10$.
			\end{itemize}
		\item SLA
			\begin{itemize}
				\item Rogers: RogersRadix -- $\varepsilon_r = 2.8$.
				\item 3Dresyns:  $\varepsilon_r = 8-9$ with $\tan \delta < 0.04$, pricey
			\end{itemize}
		\item SLS
			\begin{itemize}
				\item Ceramics are difficult to print, more suitable to binder jetting,
				\item Not considered.
			\end{itemize}
  \end{itemize}
\end{frame}

\begin{frame}[fragile]
  \frametitle{Conductive Materials}
  \begin{itemize}
    \item Carbon-loaded filaments possess low conductivity and behave primarily as lossy dielectrics.
			\begin{itemize}
				\item These low-conductivity materials excel as monolithic electromagnetic absorbers.
				\item Proto-pasta has PLA loaded with carbon black.
			\end{itemize}
    \item Copper-loaded composites achieve significantly higher bulk conductivities.
			\begin{itemize}
				\item High-conductivity materials effectively mimic perfect conductors at microwave frequencies, enabling single-step metamaterial fabrication.
				\item They are pricey, not considered.
			\end{itemize}
  \end{itemize}
\end{frame}

\begin{frame}[fragile]
  \frametitle{LWA: Monolithic Multi-Material Matched Terminations}
	\section{Leaky-wave Antennas}
  \begin{itemize}
    \item Unradiated residual power at the aperture end reflects and generates high sidelobes \footfullcite{Monticone2015}.
    \item Multi-head fused deposition modeling permits co-printing a structural dielectric and a lossy composite absorber.
    \item Continuous spatial impedance tapering suppresses back-reflection and ensures reflectionless energy transfer.
    \item Static dissipative polylactic acid or conductive polyurethane can be utilized for distributed attenuation.
  \end{itemize}
\end{frame}

\begin{frame}[fragile]
  \frametitle{LWA: Continuous Sinusoidal Reactance Surfaces}
  \begin{itemize}
		\item Most basic LWAs discretize periodic modulation into square pulses (blocks of different $\varepsilon_r$ materials, air), exciting unwanted higher-order space harmonics.
		\item Commonly a more complex unit cell design or heigh modulation in $z$-axis is used to solve this.
		\item Different Approach: Varying infill or using foaming PLA
		\item More complex solutions then already existing, but final structure is a slab on the outside
  \end{itemize}
\end{frame}

\begin{frame}[fragile]
  \frametitle{LWA: Chirped Conformal Pre-Distorted Geometries}
  \begin{itemize}
    \item Mounting a periodic leaky-wave antenna onto a curved surface introduces severe non-linear phase errors.
    \item Additive manufacturing enables phase correction by continuously modifying the internal geometric fill fraction \footfullcite{Bartley2025}.
    \item Inverse-design methodology progressively chirps the unit-cell periodicity along the curved propagation axis.
    \item The conformal structure directly synthesizes a collimated high-gain plane wave from a highly curved physical surface.
  \end{itemize}
\end{frame}

\begin{frame}[fragile]
  \frametitle{Reflectarrays: Graded-Index via Infill Density Control}
	\section{Reflectarrays Antennas}
  \begin{itemize}
    \item Conventional dielectric reflectarrays rely on geometric height variations that introduce shadowing and diffraction.
    \item Spatially modulating the internal fill fraction tunes the local effective permittivity while maintaining a flat external profile.
    \item This continuous tuning enables completely flush-mounted reflectarrays for conformal aerospace applications.
    \item The process utilizes high-permittivity ceramic filaments and requires programmatic slicer automation \footfullcite{Beccaria2025}.
  \end{itemize}
\end{frame}

\begin{frame}[fragile]
  \frametitle{Reflectarrays: Sub-THz Structures for 6G}
  \begin{itemize}
    \item Printed circuit board fabrication tolerances and conductor losses become prohibitive above 100 GHz.
    \item Spatial feeding in reflectarrays entirely eliminates the lossy corporate feed network.
    \item All-dielectric unit cells mitigate high-frequency conductor losses.
    \item This direction establishes the practical upper frequency limits of additive manufacturing for telecommunication systems.
  \end{itemize}
\end{frame}

\begin{frame}[fragile]
  \frametitle{Reflectarrays: Conformal Curved Structures}
  \begin{itemize}
    \item Standard planar phase synthesis produces severe beam defocusing when applied to arbitrarily curved surfaces.
    \item Phase distributions must be recalculated -- to account for local surface normal rotation.
			\begin{itemize}
				\item Metasurface with high phase stability is required -- range if incidence angle is larger then in planar arrays
			\end{itemize}
		\item Additive manufacturing permits direct fabrication onto pre-shaped (with significant slicer modification) substrates or as free-standing curved structures.
    \item The large design space necessitates programmatic computer-aided design and surrogate-based optimization algorithms.
  \end{itemize}
\end{frame}

\begin{frame}[fragile]
  \frametitle{Reflectarrays: Integrated Dielectric Rod Feed}
  \begin{itemize}
    \item Conventional commercial horn antennas introduce alignment uncertainty and prevent complete system miniaturization.
    \item A dielectric rod antenna can be printed simultaneously as a vertical protrusion directly from the array surface.
    \item This monolithic fabrication eliminates alignment variables and reduces overall component costs.
    \item The fixed focal length demands extreme simulation accuracy prior to the printing process.
  \end{itemize}
\end{frame}

\appendix
\begin{frame}[plain]
  \frametitle{Bibliography}
  \section{Bibliography}
  \changefontsizes{7pt}
  \printbibliography
\end{frame}

\end{document}
